\section{Chapitre~\ref{sec:dofaalt}. Autres théories de \nt{DOPA}}

Chacune des images élargies s'appuie sur ses propres mesures directes de la dopamine (dispersion des points d'inversion, réponses au désagréable et au nouveau, montée lente, réponses au changement de goût), mais les formules qui les expliquent sont des théories, et pour deux d'entre elles les expériences se contredisent encore.

{\footnotesize
\setlength{\tabcolsep}{3pt}\renewcommand{\arraystretch}{1.15}
\begin{longtable}{>{\raggedright}p{0.5cm}>{\raggedright}p{4.0cm}>{\raggedright}p{5.6cm}>{\raggedright}p{2.3cm}>{\raggedright\arraybackslash}p{1.7cm}}
\toprule
N\textsuperscript{o} & Affirmation & Expérience et ce qui est mesuré & Source & Statut \\
\midrule\endfirsthead
\toprule
N\textsuperscript{o} & Affirmation & Expérience et ce qui est mesuré & Source & Statut \\
\midrule\endhead
1 & La règle asymétrique~\eqref{eq:asymmetric} converge vers le point~\eqref{eq:expectile}; pour $\kappa=1/2$ c'est la moyenne, pour une goutte de probabilité $1/2$, $V=\kappa$ (fig.~\ref{fig:expectiles}) & Le point~\eqref{eq:expectile} est un expectile, solution d'un problème de moindres carrés avec des poids différents pour les écarts positifs et négatifs; pour l'exemple du chapitre, démonstration dans le chapitre même. & \cite{NeweyPowell1987}; démonstration dans le chapitre & théorie \\
2 & Les différents neurones \nt{DOPA} ont des points d'inversion différents, ordonnés selon l'asymétrie (proposition~\ref{prop:expectile}) & Souris, neurones \nt{DOPA} de l'aire tegmentale ventrale marqués par optogénétique, récompenses de tailles différentes: le point où le pic fait place au creux diffère selon les neurones; il est plus haut chez les neurones aux réponses plus asymétriques; les réponses du groupe permettent de reconstituer la forme de la distribution des récompenses. Que ce soit la fonction principale de la dispersion est une hypothèse. & \cite{Dabney2020} & mesuré \\
3 & Une partie des neurones \nt{DOPA} répond au désagréable par un pic & Singes: certains neurones \nt{DOPA} sont excités par l'indice d'une récompense et inhibés par l'indice d'un jet d'air dans l'œil, d'autres sont excités par les deux; les groupes se trouvent dans des parties différentes du mésencéphale. & \cite{Matsumoto2009, BrombergMartin2010} & mesuré \\
4 & Le module de l'erreur ou la nouveauté fixe la vitesse d'apprentissage & Les travaux cités ne contiennent pas d'expérience directe où le signal d'importance des neurones \nt{DOPA} change la vitesse d'apprentissage; la revue propose un rôle de signal \og attention, c'est important\fg{}. & \cite{BrombergMartin2010} & hypothèse \\
5 & Un fil séparé pour l'axe du danger & Souris: les neurones \nt{DOPA} qui projettent vers l'extrémité postérieure du striatum répondent aux stimuli nouveaux et intenses, et non à la récompense; leur destruction affaiblit l'évitement d'un objet nouveau. & \cite{Menegas2018} & mesuré \\
6 & Temps d'action optimal $\tau_a^*=\sqrt{C/\bar R}$~\eqref{eq:vigor} & Minimum de la somme $C/\tau_a+\bar R\tau_a$; démonstration dans le chapitre. Dans le modèle d'origine, la vitesse des actions est fixée par le prix du temps, égal à la récompense moyenne. & \cite{Niv2007}; démonstration dans le chapitre & théorie \\
7 & Le niveau lent de dopamine porte $\bar R$ et fixe la vitesse des actions (proposition~\ref{prop:multiplex}) & Rats, microdialyse et voltampérométrie dans le noyau accumbens: le niveau de dopamine suit, de minute en minute, la fréquence des récompenses et la vitesse des actions. Chez l'homme, la L-DOPA renforce la dépendance du temps de réaction à la fréquence moyenne des récompenses. Que le niveau lent soit précisément $\bar R$ n'est pas démontré. & \cite{Hamid2016, Beierholm2013vigor} & hypothèse \\
8 & Le niveau lent a deux sources: la libération varie aux sorties sans que change la fréquence des impulsions, mais une montée lente existe aussi dans la fréquence elle-même & Rats, même tâche: la libération dans le noyau accumbens suit l'attente changeante de la récompense, mais pas la fréquence des impulsions des neurones \nt{DOPA} identifiés de l'aire tegmentale ventrale. Souris dans un couloir virtuel (ligne~10): la montée progressive à l'approche de la récompense existe aussi dans les impulsions de neurones \nt{DOPA} identifiés. & \cite{Mohebi2019, Kim2020} & mesuré \\
9 & La dopamine monte progressivement pendant que l'animal s'approche d'une récompense lointaine & Rats dans un labyrinthe, voltampérométrie dans le striatum: la concentration de dopamine augmente en quelques secondes à mesure que le but approche; la pente dépend de la distance et de la taille de la récompense. & \cite{Howe2013ramp, Hamid2016} & mesuré \\
10 & La montée est l'erreur $\delta$ quand l'estimation se précise~\eqref{eq:ramp}; un saut dans le couloir donne un pic (fig.~\ref{fig:ramp}) & La formule et la valeur $\delta=0.75\,V$ par seconde sont démontrées dans le chapitre. Expérience: souris dans un couloir virtuel, téléportation instantanée plus près du but; les impulsions des corps cellulaires, le calcium des corps et des terminaisons et la concentration dans le striatum répondent comme une erreur, et non comme l'attente $V$. Laquelle des deux explications vaut pour toutes les sorties est en discussion. & \cite{Kim2020} & mesuré \\
11 & Regard en arrière: la dopamine signale la mesure du lien causal~\eqref{eq:retro} & Souris, libération de dopamine dans le noyau accumbens dans des expériences où cette théorie et l'erreur~\eqref{eq:ctd} prédisent des choses différentes: dans plusieurs expériences, la libération a suivi la première. & \cite{Jeong2022} & hypothèse \\
12 & Au cours de l'apprentissage, la réponse dopaminergique se déplace progressivement de la récompense vers l'indice & Souris, signal des terminaisons des neurones \nt{DOPA} dans le striatum ventral au fil de l'apprentissage: le pic passe progressivement du moment de la récompense au moment de l'indice, comme l'exige l'apprentissage selon~\eqref{eq:ctd}. & \cite{Amo2022} & mesuré \\
13 & Les neurones \nt{DOPA} répondent au changement de goût de la récompense à valeur égale & Rats, enregistrement de neurones de l'aire tegmentale ventrale: le remplacement d'un jus par un autre de même valeur provoque une réponse, bien que l'erreur~\eqref{eq:ctd} soit nulle. & \cite{Takahashi2017} & mesuré \\
14 & Des pics artificiels enseignent le lien entre deux stimuli neutres & Rats, expérience de préconditionnement sensoriel sans récompense: l'activation optogénétique des neurones \nt{DOPA} lors du passage du premier stimulus au second crée le lien, leur inactivation empêche de l'apprendre. & \cite{Sharpe2017} & mesuré \\
15 & Vecteur d'erreurs par caractéristique~\eqref{eq:vectortd} & Théorie de l'erreur de prédiction des caractéristiques; la forme vectorielle n'a pas été testée directement. & \cite{Gardner2018} & hypothèse \\
16 & Dans une même tâche, différents neurones \nt{DOPA} portent différentes grandeurs & Souris dans un couloir virtuel, imagerie à deux photons des neurones \nt{DOPA}: certains sont liés à la récompense, d'autres à la vitesse et aux virages, d'autres encore aux indices et à la justesse du choix. & \cite{Engelhard2019} & mesuré \\
17 & Un faisceau au lieu d'un fil (proposition~\ref{prop:bundle}) & Synthèse des lignes~2--16: chaque décomposition est mesurée séparément; aucune expérience ne donne une image unifiée où elles seraient toutes les parties d'un même signal. & --- & hypothèse \\
\bottomrule
\end{longtable}}
